% ATTEMPT_STATUS: unresolved
% TOP_PROBLEM: {"problemId": "problem.local-smoothing-conjecture-for-wave-equations", "canonicalTitle": "Local Smoothing Conjecture for Wave Equations", "releaseRank": 201, "primaryDomain": "analysis_pde", "primaryDomainLabel": "Analysis & PDE", "displayStatus": "open", "releaseStatus": "open"}
% !TeX program = lualatex
% Complete problem section copied verbatim from research_batch_201_202.tex.
\documentclass[11pt]{article}
\usepackage[margin=25mm]{geometry}
\usepackage{fontspec}
\setmainfont{Latin Modern Roman}
\usepackage{amsmath,amssymb,mathtools}
\usepackage{unicode-math}
\setmathfont{Latin Modern Math}
\usepackage{enumitem,needspace}
\usepackage{xurl,hyperref}
\hypersetup{colorlinks=true,urlcolor=blue}
\setcounter{secnumdepth}{1}
\setlength{\parindent}{0pt}
\setlength{\parskip}{5pt}
\setlength{\emergencystretch}{3em}
\setlist[itemize]{leftmargin=3em,itemsep=5pt,topsep=4pt,parsep=0pt}
\allowdisplaybreaks
\newcommand{\N}{\mathbb{N}}
\newcommand{\Z}{\mathbb{Z}}
\newcommand{\Q}{\mathbb{Q}}
\newcommand{\R}{\mathbb{R}}
\newcommand{\C}{\mathbb{C}}
\newcommand{\F}{\mathbb{F}}
\newcommand{\A}{\mathbb{A}}
\newcommand{\PP}{\mathbb P}
\newcommand{\E}{\mathbb{E}}
\newcommand{\eps}{\varepsilon}
\newcommand{\dd}{\,\mathrm d}
\newcommand{\sref}[2]{\hyperlink{source:#1:#2}{[S#2]}}
\newcommand{\eref}[2]{\hyperlink{extra:#1:#2}{[E#2]}}
% Turn the next switch off for a shorter reading copy. The quotations preserve
% every exactTarget field, including qualifications omitted from a short summary.
\newif\ifshowcatalogue
\showcataloguetrue
\newcommand{\cataloguescope}[1]{\ifshowcatalogue
	\paragraph{Exact scope in the supplied catalogue.}
	\begin{quote}\small #1\end{quote}\fi}
\begin{document}
\setcounter{section}{200}
	% ================================================================
	% problemId: problem.local-smoothing-conjecture-for-wave-equations
	% source releaseRank: 201; source releaseStatus: open
	% exactTarget SHA-256: 69808b9a81b2b9968148156e4e2a52f3bd538cfa32b8ff88f04e4a834f21f1b0
	\clearpage
	\section{\texorpdfstring{Local Smoothing Conjecture for Wave Equations}{Local Smoothing Conjecture for Wave Equations}}\label{201}
	\subsection{Definitions and mathematical statement}
	For a Schwartz function on $\R^n$, the half-wave propagator is the Fourier multiplier
	$\widehat{e^{it\sqrt{-\Delta}}f}(\xi)=e^{it|\xi|}\widehat f(\xi)$, using the Fourier convention in which $-\Delta$ has symbol $|\xi|^2$. Define the Bessel-potential norm
	$\|f\|_{W^{s,p}}=\|(1-\Delta)^{s/2}f\|_{L^p}$ and put $s_p=(n-1)(1/2-1/p)$. The assertion is
	\[
	\left(\int_1^2\int_{\R^n}|e^{it\sqrt{-\Delta}}f(x)|^p\dd x\dd t\right)^{1/p}
	\le C_{n,p,\eps}\|f\|_{W^{s_p-1/p+\eps,p}}
	\]
	for $n\ge3$, $p>2n/(n-1)$, and every $\eps>0$. Constants are independent of $f$. This is a spacetime estimate with an arbitrarily small derivative loss; the excluded endpoint and a zero-loss bound are not asserted.
	\par\smallskip\textit{Formulation and concept references:} \sref{201}{1}.
	\cataloguescope{For every n >= 3, p > 2n/(n-1), and epsilon > 0, prove ||exp(it sqrt(-Delta))f||\_\{L\textasciicircum{}p(R\textasciicircum{}n x [1,2])\} <= C\_\{p,epsilon\} ||f||\_\{W\textasciicircum{}\{s\_p-1/p+epsilon,p\}(R\textasciicircum{}n)\}, where s\_p=(n-1)(1/2-1/p).}
	\subsection{Short English statement}
	Does averaging the wave evolution over time save almost one-over-p derivatives throughout the conjectured range of exponents?
	% BEGIN RESEARCH ADDITION 201
	\subsection{Research attempt}
	
	\paragraph{Current frontier and source audit (24 September 2026).}
	The published cone decoupling theorem of Bourgain--Demeter gives the desired
	almost-$1/p$ gain in the range $p\geq 2(n+1)/(n-1)$; this is not the full
	range in this entry.\ \eref{201}{1}
	The more recent manuscript of Gan--He--Li--Wu, version 2 of 4 January 2026,
	states the improved range $p\geq P_n$, where
	\begin{equation}\label{ra201:frontier}
		P_n=\begin{cases}
			2+\dfrac{8}{3n-3},&n\text{ odd},\\[4pt]
			2+\dfrac{8}{3n-2-\omega(n)},&n\text{ even},
		\end{cases}
		\qquad \omega(n)=\frac{4(7n-2)}{3n^2+5n+14}.
	\end{equation}
	Its Theorem 1.3 treats both flat and compact-manifold problems. In particular,
	$P_3=10/3$, whereas the present flat-space target begins at $p>3$.
	The manuscript's sharp odd-dimensional \emph{manifold} result must not be
	read as a resolution of the stronger Euclidean range. This notebook uses
	its stated theorem as recent frontier information, not as an independently
	refereed proof supplied here.\ \eref{201}{2}
	
	Radial inputs are a genuinely easier, already studied subcase; the 2025
	paper of Beltran--Roos--Rutar--Seeger proves even a fractal-time radial
	extension.\ \eref{201}{3} The attack below rederives a radial estimate by
	an elementary one-dimensional calculation, and then tries to remove the
	radial hypothesis. That calculation is not claimed as a new theorem.
	The publisher text at the original \sref{201}{1} was not retrievable in
	this session. No missing definition is filled in from it: the displayed,
	self-contained target in the supplied document is used verbatim.
	
	\paragraph{Attack map.}
	Write $U_t=e^{it\sqrt{-\Delta}}$ and
	\[
	p_* =\frac{2n}{n-1},\qquad
	a_n(p)=s_p-\frac1p=\frac{n-1}{2}-\frac np.
	\]
	Three possible routes were considered. Interpolation from an available
	large-$p$ theorem would need a better low-$p$ input; the exact loss is
	computed below. A second route transports each narrow angular cap almost
	without distortion and then tries to exploit time averaging when the caps
	are recombined; its missing ingredient is a deterministic, phase-sensitive
	synthesis estimate, not an estimate for an individual cap. The route
	pursued most fully is to solve the radial dynamics exactly in dimension
	three and then expand in spherical harmonics. It would need uniform
	control of the angular-momentum barrier \emph{and} the summation of angular
	modes. Bounds for one fixed spherical harmonic do not supply those two
	uniformities.
	
	\paragraph{Attempt A: make the frequency and interpolation losses explicit.}
	Fix smooth dyadic annular projections $P_\lambda$, $\lambda=2^k\geq2$,
	with slightly larger annular cutoffs available. The sufficient
	frequency-localized statement is
	\begin{equation}\label{ra201:dyadic}
		\|U_t g\|_{L^p(\R^n\times[1,2])}
		\leq C_{n,p,\eta}\lambda^{a_n(p)+\eta}\|g\|_p,
		\qquad \operatorname{supp}\widehat g\subset
		\{\lambda/2\leq|\xi|\leq2\lambda\},
	\end{equation}
	for every $\eta>0$, with harmless changes of the fixed annulus. To see
	that this suffices for the \emph{stated} assertion, take
	$\eta=\eps/2$ and put $s=a_n(p)+\eps$. Annular multiplier bounds give
	\[
	\|P_{2^k}f\|_p\leq C_s2^{-ks}\|(1-\Delta)^{s/2}f\|_p.
	\]
	The triangle inequality and
	$\sum_{k\geq1}2^{-k\eps/2}<\infty$ prove the assertion. The low-frequency
	term is uniformly bounded for $1<p<\infty$ by the multiplier theorem;
	$e^{it|\xi|}$ need not be declared smooth at $\xi=0$ to use that theorem.
	One first applies the argument to finite dyadic sums and then passes in
	$L^p$ to the limit. Conversely, the requested Sobolev estimate implies
	\eqref{ra201:dyadic}, with the same arbitrarily small loss, by applying
	the Bessel multiplier on an annulus. Thus no summability problem is
	being hidden in the notation.
	
	Here is an obstruction to the first attack route. Suppose the sharp
	frequency exponent $a_n(q)$ is available at some $q>p_*$. Interpolating
	that estimate with the exact $L^2$ identity on the time interval of length
	one gives, for $2<p<q$,
	\[
	\vartheta=\frac{\frac12-\frac1p}{\frac12-\frac1q},\qquad
	b(p)=\vartheta a_n(q).
	\]
	An elementary subtraction yields
	\begin{equation}\label{ra201:interpolationloss}
		b(p)-a_n(p)=\frac{1-\vartheta}{2}>0.
	\end{equation}
	The arbitrarily small losses at $q$ do not remove this fixed positive
	loss. For example, $n=3$, $q=10/3$, $p=3$ give
	$\vartheta=5/6$ and a loss of $1/12$ derivatives. Here $p=3$ is only a
	boundary diagnostic; \eqref{ra201:interpolationloss} is also positive
	throughout the unproved open interval $3<p<10/3$. Supplying a negative
	power of $\lambda$ at $L^2$ to repair the interpolation is impossible:
	unitarity gives $\|U_tg\|_{L^2_{x,t}}=\|g\|_2$ exactly.
	
	\paragraph{Attempt B: an exact radial calculation, including its critical logarithm.}
	\textbf{Lemma 201.1 (radial benchmark, independently derived).}
	Let $n=3$, let $f$ be radial and Schwartz, and suppose its Fourier support
	lies in $\{\lambda/2\leq|\xi|\leq2\lambda\}$, $\lambda\geq2$.
	For $2\leq p<\infty$,
	\begin{equation}\label{ra201:radialbound}
		\|U_tf\|_{L^p(\R^3\times[1,2])}
		\leq C_p B_p(\lambda)\|f\|_p,
		\quad
		B_p(\lambda)=
		\begin{cases}
			1,&2\leq p<3,\\
			(1+\log\lambda)^{1/3},&p=3,\\
			\lambda^{1-3/p},&p>3.
		\end{cases}
	\end{equation}
	The implication for the target is only the radial three-dimensional
	subcase. Its relationship with the established radial theory is discussed
	in \eref{201}{3}, Theorem 1.1 and Section 5.
	
	\emph{Proof.}
	Write $f(x)=F(|x|)$ and extend $v(r)=rF(|r|)$ as an odd Schwartz
	function on $\R$. The three-dimensional radial Fourier formula shows
	that $\widehat v$ is supported where
	$\lambda/2\leq|\rho|\leq2\lambda$. Equivalently, conjugating the radial
	Laplacian by multiplication by $r$ gives the one-dimensional Dirichlet
	Laplacian, or the ordinary Laplacian on odd functions. Let $g$ be the
	positive-frequency part of $v$. Since the Fourier support avoids zero,
	this part can be obtained using a fixed smooth multiplier
	$m(\rho/\lambda)$ which is one on the positive annulus and zero on the
	negative annulus. Oddness implies $v(r)=g(r)-g(-r)$, hence exactly
	\begin{equation}\label{ra201:radialidentity}
		rU_tf(r)=g(r+t)-g(t-r),\qquad r>0.
	\end{equation}
	This is an identity for the half-wave, not a replacement of it by the
	cosine propagator.
	
	For $p\geq2$ put $W(r)=(1+|r|)^{2-p}$. The kernel of the smooth
	positive-frequency cutoff is $K_\lambda(r)=\lambda K(\lambda r)$,
	where $K$ is Schwartz. The elementary weight comparison
	\[
	\frac{W(x)^{1/p}}{W(y)^{1/p}}
	\leq (1+|x-y|)^{(p-2)/p}
	\]
	and Young's inequality prove
	\begin{equation}\label{ra201:weightedcutoff}
		\|g\|_{L^p(W\,\dd r)}+
		\lambda^{-1}\|g'\|_{L^p(W\,\dd r)}
		\leq C_p\|v\|_{L^p(W\,\dd r)}
		\leq C_p\|v\|_{L^p(|r|^{2-p}\,\dd r)}.
	\end{equation}
	Indeed the required weighted $L^1$ norm of $K_\lambda$ is uniformly
	bounded, and that of $K_\lambda'$ is $O(\lambda)$.
	The last norm to the $p$th power is
	$2\int_0^\infty |F(r)|^p r^2\dd r$, a fixed multiple of $\|f\|_p^p$.
	This argument deliberately uses the \emph{inhomogeneous} weight $W$;
	claiming a Hilbert-transform bound on $L^p(|r|^{2-p}\dd r)$ at $p\geq3$
	would be an unjustified shortcut.
	
	For $0<r\leq1/2$, all the arguments $t\pm r$ lie in $[1/2,5/2]$.
	The triangle inequality and then the fundamental theorem of calculus,
	applied separately, give
	\[
	\|g(t+r)-g(t-r)\|_{L^p_t([1,2])}
	\leq C_p\min\{1,\lambda r\}\|f\|_p.
	\]
	For $r\geq1/2$, substitute $s=t+r$ and $s=t-r$ in the two terms.
	Since $r^{2-p}\leq C_p(1+|s|)^{2-p}$ on these substitutions,
	\eqref{ra201:weightedcutoff} gives
	\[
	\int_1^2\int_{1/2}^{\infty}
	|g(t+r)-g(t-r)|^p r^{2-p}\dd r\dd t
	\leq C_p\|f\|_p^p.
	\]
	Using polar coordinates in \eqref{ra201:radialidentity}, the remaining
	factor is bounded by
	\begin{equation}\label{ra201:radialintegral}
		1+\int_0^{1/2}r^{2-p}\min\{1,(\lambda r)^p\}\dd r
		\leq C_p
		\begin{cases}
			1,&p<3,\\
			1+\log\lambda,&p=3,\\
			\lambda^{p-3},&p>3.
		\end{cases}
	\end{equation}
	For example, the part $r<\lambda^{-1}$ is
	$\lambda^{p-3}/3$; the other part is the elementary integral of
	$r^{2-p}$. Tonelli applies to the nonnegative integrands. This proves
	the lemma.\hfill$\square$
	
	In particular, the radial estimate has the desired power
	$a_3(p)=1-3/p$ for every $p>3$ without a dyadic power loss, and has only
	a logarithm at $p=3$. The dyadic summation already proved converts this
	to the radial version of the original Sobolev assertion. It does not
	prove an unrestricted zero-loss Sobolev estimate after summing all
	frequencies.
	
	\paragraph{Attempt C: remove radial symmetry without discarding its cancellation.}
	The exact identity above suggests expanding
	\[
	f(r\omega)=\sum_{\ell=0}^{\infty}\sum_{m=-\ell}^{\ell}
	F_{\ell m}(r)Y_{\ell m}(\omega),
	\qquad -\Delta_{S^2}Y_{\ell m}=\ell(\ell+1)Y_{\ell m}.
	\]
	For $w_{\ell m}(r)=rF_{\ell m}(r)$, the radial operator is now
	\begin{equation}\label{ra201:angularbarrier}
		L_\ell=-\frac{\mathrm d^2}{\mathrm dr^2}
		+\frac{\ell(\ell+1)}{r^2}
	\end{equation}
	with the regular, Friedrichs boundary condition at zero. The corresponding
	half-wave is $e^{it\sqrt{L_\ell}}$, not the translated function in
	\eqref{ra201:radialidentity}. On $r\asymp1$, modes $\ell\asymp\lambda$
	have a potential of order $\lambda^2$, comparable to the spectral energy.
	They cannot be treated as a uniformly small perturbation of $L_0$.
	Positivity of the potential does not imply pointwise domination of this
	oscillatory half-wave by the free half-wave.
	
	The missing step has two parts: a weighted spacetime estimate for
	$e^{it\sqrt{L_\ell}}$ which remains controlled through the turning region
	$r\asymp\ell/\lambda$, and an estimate for the \emph{coupled angular sum}
	with the $L^p$ norm of the original sum on the right. A fixed-$\ell$
	constant, however finite, is insufficient. There are $(M+1)^2$ modes up
	to degree $M$; a polynomial loss in $M$ cannot simply be absorbed into
	$\lambda^\eta$ when $M$ grows on the scale of $\lambda$. No uniform
	turning-region estimate or cancellation-preserving angular summation
	bound is established here. This is the precise point at which this
	extension stops.
	
	A complementary localization check helps to locate rather than remove
	that difficulty. For a cap around $e_n$ of angular width
	$O(\lambda^{-1/2})$, write $\xi_n=\lambda\rho$ and
	$\xi'=\sqrt\lambda\,\zeta$, where $\rho$ stays in a fixed compact
	subset of $(0,\infty)$ and $\zeta$ stays bounded. After removing the
	translation multiplier $e^{it\xi_n}$, the phase is exactly
	\begin{equation}\label{ra201:capphase}
		t(|\xi|-\xi_n)=
		\frac{t|\zeta|^2}{\sqrt{\rho^2+\lambda^{-1}|\zeta|^2}+\rho}.
	\end{equation}
	All derivatives of this phase times a fixed smooth cap cutoff are
	uniformly bounded for $1\leq t\leq2$. Its inverse Fourier transform in
	$(\zeta,\rho)$ has uniformly bounded $L^1$ norm, by integration by
	parts to order greater than $n$. Anisotropic dilation preserves the
	kernel's $L^1$ norm. Consequently
	\[
	\sup_{1\leq t\leq2}\|U_tf\|_p\leq C_p\|f\|_p
	\]
	for input supported in a slightly smaller such cap. This proves that
	an individual cap is controllable, but there are
	$O(\lambda^{(n-1)/2})$ caps. Summing their norms discards exactly the
	angular information the full conjecture asks us to preserve. The
	wave-packet and anisotropic-rescaling perspective is also central to
	\eref{201}{2}, Sections 1.3 and 5.1; no new decoupling inequality is
	being inferred from the elementary cap calculation.
	
	\paragraph{Stress test: focusing, fake dephasing, and quantifiers.}
	A counterexample search should first challenge the proposed exponent.
	Fix $t_0=3/2$ and a nonzero, nonnegative smooth radial $\psi$ supported
	in $1<|\eta|<2$, and define
	\[
	f_\lambda(x)=\int_{\R^n}
	e^{i\lambda(x\cdot\eta-t_0|\eta|)}\psi(\eta)\dd\eta.
	\]
	At $|x|,|t-t_0|\leq c/\lambda$, continuity of the phase and positivity
	of $\int\psi$ give $|U_tf_\lambda(x)|\geq c'>0$, uniformly in
	$\lambda$. Hence
	\[
	\|U_tf_\lambda\|_{L^p_{x,t}}\gtrsim\lambda^{-(n+1)/p}.
	\]
	For completeness, the input estimate is obtained by stationary phase in
	the angular variables near the two directions parallel to $x$, followed
	by integration by parts in the radial frequency. On
	$t_0/2<|x|<2t_0$ this gives, to arbitrarily high order,
	\[
	|f_\lambda(x)|\leq C_N\lambda^{-(n-1)/2}
	(1+\lambda\,||x|-t_0|)^{-N}+O_N(\lambda^{-N}).
	\]
	Outside that bounded annulus the phase gradient does not vanish and
	integration by parts gives integrable rapid decay, with arbitrary
	powers of $\lambda^{-1}$. Thus, for fixed finite $p$,
	\[
	\|f_\lambda\|_p\lesssim
	\lambda^{-(n-1)/2-1/p},\qquad
	\frac{\|U_tf_\lambda\|_{L^p_{x,t}}}{\|f_\lambda\|_p}
	\gtrsim\lambda^{a_n(p)}.
	\]
	The attempt to beat the proposed power fails: this is a sharpness test,
	not a counterexample to the stated estimate. For $n=3$ it agrees with
	the radial upper bound just proved.
	
	Another attempted shortcut was to treat radial-frequency pieces as if
	time averaging provided arbitrary-$L^p$ orthogonality. The proposed
	scalar ingredient
	\[
	\Big\|\sum_{j=1}^{N}a_je^{ij(t-t_0)}\Big\|_{L^p([1,2])}
	\leq C_\eta N^\eta\Big(\sum_{j=1}^{N}|a_j|^2\Big)^{1/2}
	\]
	is false for $p>2$ and $\eta<1/2-1/p$. Set every $a_j=1$:
	on $|t-t_0|\leq c/N$ the sum has modulus at least $c'N$, so the
	left side is at least $c''N^{1-1/p}$, whereas the right side is
	$C_\eta N^{1/2+\eta}$. This refutes that intermediate lemma, not
	local smoothing. A successful time-frequency argument must retain the
	spatial constraints and cannot use this scalar assertion.
	
	The radial proof controls the region $r=0$ through the difference in
	\eqref{ra201:radialidentity}, rather than estimating its two summands
	separately in a singular weight. The cap proof uses a smooth cutoff
	supported away from $\rho=0$. Constants may depend on the fixed exponent
	$p$, which is allowed; they may not depend on angular momentum or the
	number of caps in a claimed unrestricted estimate. Finally, an
	endpoint bound at $p_*$ with arbitrarily small loss would be a useful
	\emph{stronger sufficient goal}, but the supplied problem excludes that
	endpoint and it is not silently added to the target.
	
	\paragraph{Outcome.}
	\textbf{Outcome: Exploratory approach with unresolved gap.}
	
	The notebook establishes the dyadic reduction, the exact interpolation
	loss, the radial three-dimensional estimate with its critical logarithm,
	and uniform transport for one narrow angular cap. It also supplies
	explicit tests that defeat a better-than-allowed exponent and a proposed
	scalar dephasing lemma. The radial estimate belongs to an already known
	special-case theory; these calculations are not advertised as a new
	advance in the unrestricted frontier.
	
	The unresolved step is uniform angular-momentum control together with
	cancellation-preserving synthesis, or a genuinely different argument
	that supplies \eqref{ra201:dyadic} throughout
	$p_*<p<P_n$. Neither this interval for arbitrary inputs in dimension
	three nor the all-dimensional target has been proved or disproved.
	No numerical experiment has been used as a substitute for either claim.
	% END RESEARCH ADDITION 201
	\subsection{Sources}
	\begin{itemize}\raggedright
		\item[\textnormal{[S1]}]\hypertarget{source:201:1}{} Formal statement, Introduction, local smoothing conjecture.
		\url{https://www.sciencedirect.com/science/article/pii/S0022123623003786}.
		{\small\textit{Catalogue formulation source.}}
		% BEGIN RESEARCH REFERENCES 201
		\item[\textnormal{[E1]}]\hypertarget{extra:201:1}{}
		J. Bourgain and C. Demeter, \emph{The proof of the $\ell^2$ Decoupling
			Conjecture}, Annals of Mathematics (2) 182 (2015), no.~1, 351--389,
		Theorem 1.2 (cone decoupling) and the discussion immediately following it.
		\url{https://annals.math.princeton.edu/2015/182-1/p09};
		author manuscript \url{https://arxiv.org/pdf/1403.5335v3}.
		{\small\textit{Published result; consulted 24 September 2026.}}
		\item[\textnormal{[E2]}]\hypertarget{extra:201:2}{}
		S. Gan, D. He, X. Li, and S. Wu, \emph{On local smoothing estimates for
			wave equations}, arXiv:2502.05973v2, 4 January 2026, Theorem 1.3,
		equation (1.4), and Sections 1.1, 1.3, and 5.1.
		\url{https://arxiv.org/html/2502.05973v2}.
		{\small\textit{Recent manuscript theorem, not independently audited here;
				version 2 has four authors and an improved main result. Consulted
				24 September 2026.}}
		\item[\textnormal{[E3]}]\hypertarget{extra:201:3}{}
		D. Beltran, J. Roos, A. Rutar, and A. Seeger, \emph{A fractal local
			smoothing problem for the wave equation}, Bulletin of the London
		Mathematical Society 57 (2025), no.~12, 3667--3690; author manuscript
		arXiv:2501.12805v1, 22 January 2025, Theorem 1.1, equation (1.6),
		and Section 5 (radial inputs).
		\url{https://arxiv.org/html/2501.12805v1}.
		{\small\textit{The radial theorem, not its January 2025 general-frontier
				summary, is used here. Consulted 24 September 2026.}}
		% END RESEARCH REFERENCES 201
	\end{itemize}
	
\end{document}
